\chapter{Installing and running}\label{ch:install}
\begin{goals}
\textbf{In this chapter} you choose how to run wuffle (a file, the website, an installed app, or the
Android shell), learn where your work is kept, and what each browser can and can't do.
\end{goals}

wuffle is one file, \file{wuffle.html}, about 770~KB. Everything it needs is inside it: the program, the
World Magnetic Model, the Tocher sample, the shell. wuffle itself fetches nothing when it opens
(\cref{sec:local}); the website's copy has a service worker that checks for a newer version in the
background (\cref{sec:run-site}). There is no installer, so \enquote{installing} wuffle means choosing
where the file comes from.

\section{Three ways to run it}\label{sec:run-ways}

\subsection{As a file}

Save \file{wuffle.html} anywhere and open it: double-click it, or drag it onto a browser window. It runs
offline from the first moment, because it never needed the network. You can copy it to a USB stick,
attach it to an email, or keep one beside each project; every copy is the whole program.

A file is also the simplest way to keep a version: the copy you used for a report still opens, years
later, exactly as it was.

\subsection{From the website}\label{sec:run-site}

The same file is published at \file{gentropic.org/wuffle}. Opened there, it behaves as the file does,
with two additions:

\begin{description}
\item[It keeps working offline.] On the first visit the browser stores the page (a \term{service
worker} does this; it is part of the site, not of wuffle). After that, the address opens with no
network, in the field too.
\item[It can be installed as an app.]\index{installing as an app} Browsers that install web apps offer
it for this one: in Chrome and Edge on a computer, an install button at the right of the address bar or
an install item in the browser's menu (its wording changes between versions); in Chrome on Android,
\ui{Add to Home screen} or \ui{Install app} in the menu. The installed wuffle opens in a window of its
own, with its own icon, and still works offline.
\end{description}

\begin{note}
Installing a web app copies nothing new to your machine beyond what the browser already stored on the
first visit. It changes how wuffle opens, not what it is. Uninstall it from the same browser menu, or
from the system's app list on Android.
\end{note}

\subsection{Inside the lead-acid shell}\label{sec:run-shell}

On Android, wuffle also runs inside \term{lead-acid}, a small Android app that hosts single-file GCU
instruments and lends them what a browser can't. wuffle notices the shell and uses it. You work the same
way, but the compass gains:

\begin{itemize}
\item \textbf{The fused orientation sensor}, read directly, with the magnetometer's own calibration
report: \ui{compass low — wave the phone in a figure-8 to calibrate} when it needs it
(\cref{sec:measuring}).
\item \textbf{A position} from the phone's GNSS, stamped on every measurement, and with it \textbf{true
north} from the World Magnetic Model (\cref{sec:true-north}).
\item \textbf{A signature}: the published log is signed with a key in the phone's secure hardware
(\cref{sec:publish}).
\item \textbf{Publishing to the phone's own folders}: the log and its signature to \file{Downloads},
the photos to \file{Pictures}, where a file manager or a cable finds them.
\item \textbf{Opening shared files}: a CSV shared to wuffle, or opened with it from a file manager, is
read back into the log.
\item \textbf{The screen kept on} while the compass is open.
\end{itemize}

Outside the shell, in an ordinary phone browser, the compass still measures, but without a position (so
readings stay magnetic), without a signature, and publishing downloads the files instead.

\section{Updating}\label{sec:updating}\index{updating}

\begin{description}
\item[A file] is updated by replacing it with a newer \file{wuffle.html}. Your projects are files of
their own (\cref{sec:project-file}), so nothing is lost; the kept copy of your last session
(\cref{sec:autosave}) belongs to the place the file opens from, so it carries over when the new file
takes the old one's place.
\item[The website] updates itself. Each time you open it with a network, the browser shows the copy it
already has, at once, and checks for a newer one in the background. When there is one, it is what opens
the \emph{next} time. There is no prompt and nothing to click.
\end{description}

\begin{tip}
To use a new version straight away, open the site, wait a few seconds, and reload it.
\end{tip}

The foot of the compass page names the build\index{version and build}, as in
\ui{wuffle · 0.1.0 · d772489 · 2026-11-26}: the version, a fingerprint of the build, and its date. On
the desk, \menu{Help>About wuffle} gives the same line under its description.

\section{Where your work is kept}\label{sec:where-kept}

Three places, from the most fleeting to the most durable:

\begin{description}
\item[The tab.] What you see is held in the page while it is open.
\item[The kept copy.] A second after each change, the desk copies the whole project into the browser's
storage (\cref{sec:autosave}). Close the tab and reopen wuffle from the same place, and your session is
back. Each \file{wuffle.html} on your disk keeps its own copy, under its own path, so a copy beside each
project keeps that project's session; the website keeps one copy for the site; and a page that carries a
project (\cref{sec:save-page}) keeps a copy of its own, so two such pages never overwrite each other's
work.
\item[Your files.] A \file{.zip.wuffle} project, a \file{.wuffle.html} page, an exported figure: these
are yours, in your folders, and nothing in the browser can take them away.
\end{description}

\begin{caution}
The kept copy belongs to the browser and to where the page opened from. A file on your disk and the
website keep separate copies, and so do two copies of the file in different folders; moving or renaming
the file starts a new one. A private window keeps none past its closing, and clearing the browser's data
deletes it. It is a convenience between saves: save the project for anything you mean to keep.
\end{caution}

\subsection{Moving to another computer}

Save the project (\menu{File>Save project (.zip.wuffle)}) and open the \file{.zip.wuffle} on the other
computer, in any copy of wuffle. Or save it as a page (\menu{File>Save as page (one .html, the project
inside)}) and take that single file: it is wuffle and the project together, and opens straight into the
data.

The kept copy does not travel. Neither does the undo history: a project opened anywhere starts with a
fresh one (\cref{sec:undo}).

\section{Browsers}\label{sec:browsers}

\subsection{The desk}

The desk runs in any current browser: Chrome, Edge, Firefox or Safari. It needs one that can unpack
gzip by itself (the program inside the file is compressed) and is recent enough for the rest of it:
Chrome and Edge 93 or later, Firefox 113 or later, Safari 16.4 or later. In an older browser the page stays blank.

\subsection{The compass}\label{sec:compass-browsers}

The compass needs the phone's orientation sensors, and browsers differ more there:

\begin{description}
\item[Chrome on Android] is what the compass is made and tested on. It reads the
compass-referenced orientation the browser reports (\code{deviceorientationabsolute}); inside the
lead-acid shell, the fused sensor directly.
\item[iPhone and iPad.] Safari only reports orientation to a page that asks the user's permission first,
and wuffle doesn't ask yet. The compass reads nothing there. Typed measurements
(\cref{sec:manual-entry}) and the whole desk work.
\item[Other Android browsers] may report only orientation relative to where the phone pointed when the
page opened, not to north. When that is all wuffle gets, it says so once (\ui{this browser's compass
has no north — dips and plunges are good, azimuths are not; measurements are marked north = none}) and
marks each reading it takes with \col{north} = \code{none}. Use Chrome, or the shell.
\end{description}

Browsers give orientation only to pages they consider secure: the site, served over https, is. A copy
served from a plain \file{http://} address isn't, and its compass stays blank.

\begin{why}
wuffle could ask for the iPhone's permission, but it would then also need to say when the reading is
relative, not absolute, and it doesn't yet know how to on that platform. A compass that might point
anywhere is worse than none, so it shows nothing rather than a number that looks right.
\end{why}
